% Status: appendix_or_discussion_ready (Monitor/Swiss) + Narrativ-Entscheide.
\chapter{Erweiterungen: Monitor, Schweizer Abstimmung und praktischer Beitrag}
\label{ch:erweiterungen}

Dieses Kapitel ordnet vier Erweiterungen ausserhalb des empirischen Kerns ein.
Sie stützen die Forschungsfrage, ersetzen aber nicht den H1--H3-Beleg.

\section{Schweizer Referendum als begrenzte Fallstudie}
Als aktueller, eigenständiger Vergleich dient die eidgenössische Abstimmung vom
14.~Juni 2026. Beim offiziellen Ja-Anteil von 45.21\% lag die finale Umfrage als
Stimmenanteil näher am Resultat, während Polymarket im lokalen Live-Fenster (Ja
zuletzt 21.5\%) als binäre Ablehnungswahrscheinlichkeit klarer auf der richtigen
Seite lag \citep{zotero_poly_002}. Wichtig: Polymarket-Preise sind
Annahmewahrscheinlichkeiten, Umfragen sind Stimmenanteile. Der Binärvergleich
ist nur ein transparenter Proxy und kein Mispricing-, Effizienz- oder
Tradeability-Beweis (Abbildung~\ref{fig:swiss}).

\begin{figure}[ht]
\centering
\includegraphics[width=0.82\textwidth]{swiss_referendum_10mio_final_case_study.png}
\caption{Schweizer Referendum 14.~Juni 2026: offizielles Resultat gegen
Polymarket-Preise und Umfrage-Proxy (begrenzte Post-Resultat-Fallstudie).
Quelle: \texttt{data/results/swiss\_referendum\_10mio\_final\_case\_study.png}.}
\label{fig:swiss}
\end{figure}

\section{Anomalie-Review-Queue als Prüf-Workflow}
Die deterministische Anomalie-Diagnostik aus Abschnitt~\ref{sec:anomalie} liefert
Hinweise, die eine nachgelagerte Review-Schicht einordnet. Die aktuelle
Review-Queue enthält drei Fälle: einen mit hoher Priorität, der als
Prüf-Kandidat markiert ist, und zwei mit mittlerer Priorität als Beobachtungs-Hinweis.
Alle drei stehen auf Quellenprüfung ausstehend. Bei der Überlappung mit den
kuratierten Referenzfällen ergeben sich ein Treffer, eine teilweise und eine
fehlende Überlappung.

Diese Queue ist ausdrücklich ein menschlicher Prüf-Workflow über begrenzte,
bereits berechnete Monitor-Artefakte. Die Risiko-Etiketten sind Beobachtungs-Marker,
keine Befunde. Die Queue beweist weder private Information noch Fehlverhalten,
Kausalität, Handelbarkeit, Profitabilität oder künftige Performance.

\section{Politik-/Geo-Monitor als Read-only-Prototyp}
Der Monitor kombiniert Marktbewegung, aggregierte Wallet-Tier-Aktivität,
Konzentration und Ereigniskontext zu deterministischen Prüf-Hinweisen
\citep{zotero_poly_006,zotero_poly_009}. Er bleibt read-only, erzeugt keine
Order- oder Trading-Pfade und liefert Prüf-Hinweise, keine kausale oder
Effizienz-Evidenz. Im Haupttext erscheint er nur als Workflow, Fälle bleiben bis
zur manuellen Prüfung im Anhang.

\section{Agenten-gestützte Pipeline als dokumentierter Vertrag}
Die ursprünglich angestrebte Multiagenten-Pipeline ist hier als bewusst
begrenzter, dokumentierter Vertrag verankert. Künftige Agenten dürfen Hinweise
sichten und einordnen, aber keine Kennzahlen berechnen. Der Vertrag definiert
vier Rollen: einen EventScout, der Ereigniskandidaten sammelt, einen
CaseNarrative, der Fälle aus begrenzten Zusammenfassungen beschreibt, einen
SkepticReviewer, der gegenprüft, und einen Orchestrator, der entscheidet. Eine
spätere, ebenfalls vertraglich begrenzte MCP-Schicht stellt nur lesende Werkzeuge
bereit (Anomalie-Review-Summary, Einzelfall, Artefaktliste, Methodengrenzen), mit
maximal 50 Zeilen je Abfrage, ohne rohes SQL, ohne standardmässige
Wallet-Adress-Ausgabe und mit Pflicht zur Protokollierung jeder LLM-Interaktion.
Beide Verträge sind als nicht implementiert markiert und bleiben Future Work, bis
der deterministische Kern und die Audit-Grenzen vollständig stehen.

\section{Praktischer Beitrag: Monitoring-Werkzeug und Bau-Methodik}
Während der Untersuchung wurde Bedarf an einem praktischen Analyse-Werkzeug für
Prognosemärkte erkennbar. Als Gestaltungsbeitrag wurde daraus ein laufendes,
read-only Webprodukt entwickelt, das die Effizienz-Perspektive der Arbeit
operationalisiert. Die Umsetzung selbst nutzte eine Multiagenten-Orchestrierung
als Entwicklungsmethode, dokumentiert wird diese als Prozessbeitrag. Die Agenten
erzeugten Software, berechneten aber keine thesis-relevanten Kennzahlen. Beide
Punkte sind Sekundärbeiträge und dürfen die Forschungsfrage nicht überlagern.
Umfang und Einordnung sind mit der Betreuung abzustimmen.
